% chapters/chapter3/3_6_suggestions.tex
% !TeX root=../../main.tex

\section{پیشنهادها برای تحقیقات آتی}

با توجه به دستاوردهای این پژوهش در توسعه معماری کنترل یکپارچه برای سیستم‌های جرثقیل سقفی، مسیرهای زیر جهت گسترش و ارتقای این سامانه در مطالعات آینده پیشنهاد می‌شود:

\begin{itemize}
    \item \textbf{پیاده‌سازی آزمایشگاهی و ارزیابی تجربی:} استقرار الگوریتم \lr{NMPC} بر روی یک نمونه فیزیکی برای سنجش پایداری سیستم در برابر نویز حسگرها، تأخیرهای ارتباطی شبکه و عدم‌قطعیت‌های مدل‌سازی نظیر اصطکاک غیرخطی ریل و مفاصل.
    \item \textbf{تنظیم تطبیقی وزن‌های بهینه‌ساز مبتنی بر مسافت:} طراحی مکانیزم زمان‌بندی وزن‌ها (\lr{Weight Scheduling}) و یا بهره‌گیری از یک جدول جستجو (\lr{Lookup Table}) جهت تنظیم برخط ماتریس‌های $Q$ و $R_{\Delta}$ متناسب با طول مسیر. این رویکرد می‌تواند مصالحه مهندسی میان دقتِ ردیابی در مسافت‌های کوتاه (نظیر $0.5$ متر) و تضمین پایداری در مسافت‌های طولانی (نظیر $5$ متر) را به‌طور کامل برطرف سازد.
    \item \textbf{ارتقای مدل به آونگ دوگانه (\lr{Double Pendulum}):} در بسیاری از کاربردهای سنگین صنعتی (مانند انتقال کانتینرها)، توزیع جرم قلاب و بار باعث بروز دینامیک آونگ دوگانه می‌شود. بسط معادلات غیرخطی و تنظیم مجدد تابع هزینه برای مهار نوسانات در دو فرکانس طبیعی متفاوت، گامی حیاتی در صنعتی‌سازی این الگوریتم است.
    \item \textbf{ادغام الگوریتم‌های پرهیز از مانع (\lr{Obstacle Avoidance}):} بهره‌گیری از سامانه‌های بینایی ماشین برای تشخیص موانع متحرک در محیط کارگاه و تزریق مستقیم این داده‌ها به عنوان قیود متغیر با زمان در مسئله بهینه‌سازی، تا سیستم بتواند مسیر ایمن را به صورت برخط بازطراحی کند.
    \item \textbf{کاهش بار محاسباتی با رویکردهای یادگیری ماشین:} استفاده از شبکه‌های عصبی عمیق برای تقریب قانون کنترل صریح (\lr{Explicit MPC}) یا به‌کارگیری الگوریتم‌های یادگیری تقویتی (\lr{Reinforcement Learning}) به منظور استنتاج سریع سیگنال‌های کنترلی، که امکان اجرای این معماری را بر روی سخت‌افزارهای صنعتی با توان پردازشی محدود فراهم می‌سازد.
\end{itemize}